% CV for the website (no references section).
% Edit this file on GitHub; the PDF (cv/cv.pdf) is rebuilt automatically.
\documentclass[11pt]{article}
\usepackage[letterpaper,margin=0.9in]{geometry}
\usepackage[T1]{fontenc}
\usepackage{mathpazo}
\usepackage{amsfonts,amsmath,amssymb}
\usepackage{microtype}
\usepackage{enumitem}
\usepackage{titlesec}
\usepackage{tabularx}
\usepackage[colorlinks=true, urlcolor=blue, linkcolor=black, citecolor=black]{hyperref}
\usepackage{xcolor}

\pagestyle{plain}
\setlength{\parindent}{0pt}
\setlist{nosep,leftmargin=1.2em}

\titleformat{\section}{\large\bfseries\scshape}{}{0pt}{}[\vspace{-0.6ex}\rule{\linewidth}{0.4pt}]
\titlespacing*{\section}{0pt}{2.2ex}{1.2ex}

% \entry{left}{right}  -- one line, right-aligned date
\newcommand{\entry}[2]{\noindent #1 \hfill #2\par}
% \course{number and title}{terms}
\newcommand{\course}[2]{\noindent\hspace*{1.2em}#1 \hfill #2\par}

\begin{document}

\begin{center}
  {\LARGE\bfseries Amer Rawashdeh}\\[0.6ex]
  Department of Mathematics, Wayne State University, Detroit, MI 48202\\
  \href{mailto:rawashdeh@wayne.edu}{rawashdeh@wayne.edu} \;|\;
    \href{mailto:amerry2015@gmail.com}{amerry2015@gmail.com} \;|\;
  \href{https://clasprofiles.wayne.edu/profile/hf1056}{Wayne State Profile} \;|\;
  ORCID: \href{https://orcid.org/0009-0006-2937-2934}{0009-0006-2937-2934}
\end{center}

\section{Research Interests}
Nonlinear partial differential equations; liquid crystal flows; nonlinear analysis;
variational and geometric methods; PDE models from continuum mechanics.

\section{Education}
\entry{\textbf{Ph.D. in Mathematics}, Wayne State University, Detroit, MI}{Expected May 2027}
\hspace*{1.2em}Advisors: Fernando Charro and Tao Huang. GPA 3.91/4.00\\
\hspace*{1.2em}Dissertation: One-dimensional models for the general Oseen--Frank energy:
existence of weak and strong solutions

\medskip
\entry{\textbf{M.S. in Mathematics}, Central Michigan University, Mount Pleasant, MI}{May 2021}
\hspace*{1.2em}Advisor: Dmitry Zakharov. GPA 3.8/4.0\\
\hspace*{1.2em}Plan B Project: Theory of valuations (after Neukirch, \emph{Algebraic Number Theory}, Ch.~II)

\medskip
\entry{\textbf{M.S. in Mathematics}, Yarmouk University, Irbid, Jordan}{September 2015}
\hspace*{1.2em}Advisor: Mowaffaq Hajja\\
\hspace*{1.2em}Thesis: Number-theoretic theorems and conjectures in rings other than the integers,
with emphasis on an analogue of Goldbach's conjecture in $\mathbb{Z}_n$

\medskip
\entry{\textbf{B.S. in Mathematics}, Yarmouk University, Irbid, Jordan}{August 2010}

\section{Publications}
\begin{enumerate}[leftmargin=1.6em]
  \item T.~Huang, B.~Rababah, and A.~Rawashdeh, On existence of chemotaxis compressible Navier--Stokes
        equations modeling vascular network formation, \emph{Journal of Differential Equations}
        \textbf{480} (2026), Article 114655. \href{https://doi.org/10.1016/j.jde.2026.114655}{doi:10.1016/j.jde.2026.114655}
\end{enumerate}

\section{Submitted Manuscripts}
\begin{enumerate}[leftmargin=1.6em]
  \item T.~Huang, B.~Rababah, and A.~Rawashdeh, Compactness of the Ericksen--Leslie system modeling
        nematic liquid crystal flows in dimension two. Submitted September 2026.
  \item T.~Huang and A.~Rawashdeh, Well-posedness of biaxial nematic flows with rough initial data.
        Submitted August 2026.
  \item F.~Charro and A.~Rawashdeh, A characterization of the determinant in terms of elementary
        symmetric polynomials. Submitted June 2026.
\end{enumerate}

\section{Talks}
\begin{enumerate}[leftmargin=1.6em]
  \item On existence of chemotaxis compressible Navier--Stokes equations modeling vascular network
        formation. SIAM Great Lakes Section Annual Meeting, Purdue University, September 2026.
  \item Tensor products in algebra and in PDE, and the Einstein notation. Informal PDE Seminar,
        Wayne State University, May 2025.
  \item Hausdorff measure and dimension, and the area and coarea formulas. Special Seminar,
        Wayne State University, May 2024.
\end{enumerate}

\section{Teaching Experience}
\entry{\textbf{Wayne State University}, Detroit, MI}{2022--present}
\textit{Graduate Teaching Assistant; Instructor of Record for each course listed}\par
\smallskip
\course{MAT 1010, Beginning Algebra}{Fall 2026}
\course{MAT 2020, Calculus II}{Spring/Summer 2026; Winter 2026; Spring/Summer 2023}
\course{MAT 2150, Differential Equations and Matrix Algebra
  }{Fall 2025}
\course{MAT 2250, Elementary Linear Algebra}{Spring/Summer 2025; Winter 2025}
\course{MAT 2010, Calculus I (coordinated flipped-classroom format)}{Fall 2024}
\course{STA 2210, Probability and Statistics}{Spring/Summer 2024}
\course{MAT 2030, Calculus III}{Winter 2024}
\course{MAT 0993, Beginning Algebra (coordinated lab-based format)}{Fall 2023}
\course{MAT 1800, Elementary Functions}{Winter 2023}
\smallskip
\begin{itemize}
  \item Tutor, Mathematics Resource Center and Mathematics Computer Lab.
\end{itemize}

\medskip
\entry{\textbf{Central Michigan University}, Mount Pleasant, MI}{2018--2021}
\textit{Teaching Assistant and Research Assistant}
\begin{itemize}
  \item Instructor of record, MTH 105, College Algebra (Fall 2020, Spring 2021).
  \item Tutor, Mathematics Assistance Center; Outstanding Tutoring Award (2021--2022).
\end{itemize}

\medskip
\entry{\textbf{Ministry of Education}, Irbid, Jordan}{2010--2018}
\textit{Mathematics Teacher}
\begin{itemize}
  \item Hitteen Secondary School (2015--2018): Grade 12 mathematics including calculus;
        Grade 11 mathematics including precalculus.
  \item Kufr-Asad Secondary School (2010--2015): Grades 8--11 algebra, analytic and synthetic geometry,
        functions, introductory statistics, and probability.
\end{itemize}

\section{Curriculum and Instructional Development}
\begin{itemize}
  \item \textbf{Calculus I flipped-classroom course}, Fall 2024. Taught a section in the first semester of the
        department's coordinated flipped Calculus~I; attended weekly course-team meetings and submitted
        written recommendations on course materials to the department.
  \item \textbf{Analysis Preliminary Exam preparation}, Fall 2023. Co-developed solutions and study materials
        for graduate students preparing for the Analysis Preliminary Examination.
  \item \textbf{Tawjihi Mathematics Study Guide}, 2012. Wrote a study guide for the mathematics component of
        Jordan's national secondary-school examination.
\end{itemize}

\section{Honors and Awards}
\entry{William Martin Borgman Endowed Scholarship for Mathematics, Wayne State University}{2026--2027}
\entry{M.\,F.~Janowitz Scholarship, Wayne State University}{2023--2024}
\entry{Outstanding Tutoring Award, Graduate Student, Central Michigan University}{2021--2022}

\section{Service}
\entry{Co-organizer, AMS Special Session on Recent Trends in Nonlinear PDEs and Applications,}{April 2027}
\hspace*{1.2em}AMS Spring Eastern Sectional Meeting\par
\entry{Teaching Subcommittee, Mathematics Strategic Planning Committee, Wayne State University}{Winter 2026}
\entry{President, SIAM Student Chapter, Wayne State University}{Fall 2025--Winter 2026}
\entry{Organizer, Informal PDE Seminar, Wayne State University}{Spring/Summer 2025}
\entry{Student Representative, Department Chair Review Committee, Wayne State University}{2024--2025}
\entry{Participant, Inclusive Excellence Listening Session, Wayne State University}{Winter 2026}
\entry{Conference support, SIAM Great Lakes Section Meeting, Wayne State University}{September 2022}

\section{Professional Memberships}
\entry{American Mathematical Society (AMS)}{2022--present}
\entry{Society for Industrial and Applied Mathematics (SIAM)}{2024--present}

\section{Languages}
Arabic (native); English (fluent)

\end{document}